\section{Related work}
\label{sec:related}

This targeted review covers research on hiring audits, representation editing, information flow, and evidence verification. We prioritise primary publications that define a comparable intervention or interpret one of our measurements. Two directly relevant public preprints are identified as such because they bear on the scope and priority of the contribution. This is a narrative review of methods and assumptions. It does not estimate a pooled effect across hiring studies.

\subsection{Résumé correspondence studies and language-model audits}

\citet{bertrand2004} varied names on otherwise comparable résumés and measured employer callbacks. \citet{kline2022} extended the correspondence design to large U.S.\ employers and documented substantial variation across firms. Both designs make the name cue experimentally useful while leaving the meaning of that cue dependent on the context. Names can carry associations with several aspects of identity, and other résumé content can carry identity cues independently. \citet{kang2016}, for example, studied how applicants altered multiple parts of their résumés to obscure racial cues. \citet{behaghel2015} found that anonymous résumés changed interview and hiring outcomes in their field experiment, with adverse effects for minority candidates among participating firms. Removing a field thus changes the evidence available to a real decision process. Its impact on fairness depends on the decision process that receives the remaining evidence.

The same identification strategy has been applied to language-model screening. \citet{armstrong2024} varied names in GPT-based screening prompts and examined ratings and recommendations. \citet{wilson2024} audited résumé retrieval through embedding similarity, examining race, gender and their intersection. These studies help explain why the model, task and endpoint matter. A retrieval rank can move even when a binary decision at one cutoff does not. A mean score can remain stable while applicants exchange ranks. Our paired outcomes therefore include continuous logits, binary choices and deterministic report bytes. We keep the comparison within each underlying résumé. Changing the name string is an intervention on one cue.

\citet{wolfe2021} showed that name frequency and tokenisation matter when contextual language models represent names. \citet{an2023} explicitly studied how demographic attributes and tokenisation length affect first-name bias estimates. That finding is directly relevant to an inference guarantee. Substituting a longer name can change both the token IDs near a field boundary and the positions used for later evidence. An audit can treat those as nuisance variables. An exact intervention has to specify how they are handled. We report token admission separately and require a common retained sequence for any comparison covered by the proof.

\subsection{Individual invariance, group outcomes and institutional context}

Fairness criteria answer different questions. \citet{hardt2016} define outcomes conditional on the true label for equalised odds and equality of opportunity. Those definitions require credible task labels and a group variable. In the causal framework of \citet{kusner2017}, counterfactual fairness asks how a prediction would change in a world with a different value of a protected attribute. Our intervention keeps every retained token fixed and alters a declared name field. This is an exact statement about the model interface and cannot determine which changes to qualifications, institutions or language would accompany any broader counterfactual. Neither of our datasets provides demographic labels or validated hiring outcomes.

\citet{raghavan2020} connect these measurement choices to the claims vendors make about pre-employment assessments. \citet{selbst2019} explain how narrowing a sociotechnical problem to a model can shift scrutiny away from the institutional processes that generate its inputs and act on its outputs. In our setting, the reviewer who marks a span and the human who handles referrals both influence an applicant's experience. Access to correction, the cost of review, and the distribution of rejected inputs are therefore part of a useful deployment evaluation.

\subsection{Learned invariance and representation editing}

Counterfactual text augmentation constructs paired training examples \citep{zmigrod2019}. \citet{garg2019} study training and evaluation for predictions robust to counterfactual text edits. Paired training can encourage equal decisions on many examples. Its effect has to be tested on edits beyond the training distribution. Parameter-efficient techniques such as LoRA \citep{hu2022}, QLoRA \citep{dettmers2023} and orthogonal fine-tuning \citep{qiu2023} change how a model is adapted. Their adaptation constraints do not themselves specify which of two named prompts must have identical outputs. An earlier subtractive-adapter experiment in our project reduced average logit variation while leaving differences in all 48 tested contexts. The present method uses frozen weights and enforces equality by the inputs supplied to the protected computation.

\citet{ravfogel2020} use iterative nullspace projection to suppress protected information available to linear classifiers. \citet{belrose2023} give an exact linear concept-erasure criterion in closed form. \citet{elazar2021} use removal interventions to investigate whether encoded information affects behaviour. These are useful distinctions for judging a fairness intervention. A good probe can find a signal that a model does not use. A failed probe may have insufficient capacity. A scalar distribution uniform over \(\{-1,1\}\) and one uniform over \(\{-2,2\}\) have the same mean, but the magnitude identifies which distribution supplied any observation. A mean-direction diagnostic therefore gives a narrower claim than an intervention which makes the observed states equal on every admitted pair.

\citet{meade2022} evaluated several debiasing methods and found that gains varied with the bias being measured and could accompany reductions in language-model performance. The public preprint by \citet{karvonen2025} is especially close in application. It evaluates models under more realistic hiring prompts and edits affine concept directions at inference, including experiments where affiliations carry a demographic cue. That approach targets a wider set of signals and reports empirical bias reduction. Our mechanism targets a declared token channel and tests exact equality on that channel. The available experiments have different populations, model settings and endpoints.

\subsection{Information flow, positions and exact baselines}

\citet{goguen1982} formalise noninterference by comparing executions that agree on the information an observer may see. \citet{sabelfeld2003} review how that principle is enforced in programs, and \citet{clarkson2010} formalise properties that refer to sets of execution traces. This is the appropriate model for our claim because a single run cannot establish name-field noninterference. We must compare two runs or prove that their protected inputs and subsequent execution coincide.

Transformer attention \citep{vaswani2017} can mix information across tokens. Rotary position embeddings assign position-dependent rotations to queries and keys \citep{su2024}. Limiting which keys remain visible while retaining original positions can still change the relative distance across a name field. Token pruning methods such as those of \citet{kim2022} also operate on a selected sequence, but choose tokens to preserve task behaviour or improve efficiency. In our setting, token selection follows a reviewed evidence policy and logical positions form part of the proof.

The nearest public mechanism is the serving-time span veto in \citet[Section III]{russinovich2026}, which verifies exact outside-span activations under a fixed-position, same-token-length substitution across multiple model families. It establishes a direct antecedent for exact marked-span isolation. We therefore locate the distinct contribution more narrowly. We define and test a name-field contract that admits unequal name lengths when the same ordered evidence IDs can be preserved, pair it with compact logical execution, and evaluate admission, score equality and evidence quality on the same interface. We include a native reference given identical retained IDs, because that reference should implement the same protected computation. Table~\ref{tab:methods} compares the evidence yielded by selected peer-reviewed methods.

\begin{table}[tbp]
\small
\centering
\caption[Comparison with published methods]{Evidence from selected published methods and the corresponding claim tested here.}
\label{tab:methods}
\begin{tabularx}{\linewidth}{@{}>{\raggedright\arraybackslash}p{.23\linewidth}
 >{\raggedright\arraybackslash}X >{\raggedright\arraybackslash}X@{}}
\toprule
Approach & What the evaluation establishes & What our experiment tests separately \\
\midrule
Correspondence and retrieval audits \citep{bertrand2004,wilson2024}
 & Sensitivity to a manipulated cue in a specified population or ranking task.
 & Exact equality of each admitted pair, including logits and aligned states.\\[2pt]
Counterfactual text training \citep{garg2019}
 & Robustness to evaluated token edits after training.
 & Equality by a deterministic inference contract, with no task-specific weight update.\\[2pt]
Linear representation erasure \citep{ravfogel2020,belrose2023}
 & Suppression of a specified class of decoders or linear concept information.
 & Equality of the full aligned state for each admitted pair after the embedding boundary.\\[2pt]
Question-answering verification \citep{wang2020,fabbri2022}
 & Agreement between a claim and an answer extracted from a document.
 & Whether exact agreement with a supplied span predicts human support labels.\\[2pt]
Selective classification \citep{geifman2017,jones2021}
 & Error as a function of the proportion of cases on which a classifier acts.
 & Separate counts for token admission, confirmation and evidence error.\\
\bottomrule
\end{tabularx}
\end{table}


\subsection{Source verification and selective classification}

SQuAD~2.0 introduced adversarially written unanswerable questions paired with plausible answers \citep{rajpurkar2018}. FEVER separates a claim's verification label from the evidence retrieved in support of that label \citep{thorne2018}. Work on attributed generation likewise measures whether a citation supports its associated statement \citep{gao2023}. These designs arise from a common problem. Locating a plausible span can be easier than judging whether the span answers the question. \citet{wang2020} and \citet{fabbri2022} use question answering to check consistency between a document and a proposed summary. Our candidate-blind verifier tests a restricted version of that idea. It must reproduce an exact proposed answer and verify the source offsets. The human-labelled benchmark then tests the semantic step independently.

Selective classification makes the risk--coverage trade-off explicit \citep{geifman2017}. \citet{jones2021} show that abstention can magnify group disparities. \citet{lee2021} study a group-sensitive selective prediction criterion. \citet{angelopoulos2025} address statistical risk calibration for predictive systems on new examples under their stated assumptions. Our statistical question is different. We sample without replacement from a frozen, finite set of confirmations and verify the audit against its entire labelled census. This avoids relying on an independence assumption about errors within one article. We do not treat four repeated names as four independent confirmations.
